\ansclass{10}{Reliability \& Resilience}
\ans{1} (a) $720/722 \approx$ \final{99.72\%}. (b) A: $1\,440/1\,442 \approx$ \final{99.86\%}; B: $720/720.017 \approx$ \final{99.998\%}. \final{Fund B first}: faster recovery buys more than fewer failures, and it is cheaper. (c) $0.0001 \times 43\,830$ min $\approx$ \final{4.4 min per month}: one manual intervention takes longer, so recovery must be \final{automatic}.
\careful{quoting nines without a window: 99.9\% is 43.8 min per month but 40.3 min per 28 days.}
\ans{2} (a) $0.001 \times 28 \times 10^6 =$ \final{28\,000 failed requests}, or $0.001 \times 28 \times 1\,440 \approx$ \final{40.3 min} of full outage. (b) $1\,000\,000/24 \approx 41\,667$ requests/h; $0.05 \times 3 \times 41\,667 \approx$ \final{6\,250 failures, $\approx$ 22\% of the budget}. (c) The SLA carries money (credits), so it promises less than the team aims for: \final{the SLO is the early warning that keeps the SLA safe}.
\careful{measuring the SLI on the health-check endpoint, or with an average, instead of on the requests fans make.}
\ans{3} (a) Burn $=$ error rate$/0.1\%$: \final{1, 6 and 14.4}; the budget lasts \final{28, $\approx$ 4.7 and $\approx$ 1.9 days}. (b) $14.4 \times 1/672 \approx$ \final{2\% of the budget} (exactly 2\% on a 30-day window). (c) The 5-minute window confirms that it is \final{still burning}, so the alert stops soon after the problem is fixed; CPU is a cause, not a symptom: \final{page on what users feel}. (d) 70\% left with a flag: \final{ship}. 8\% left before an on-sale: \final{wait} (the budget policy freezes the purchase path).
\careful{mixing windows: a burn rate computed against 30 days and a budget of 28 days.}
\ans{4} (a) $0.9999 \times 0.995^3 \times 0.999^2 \approx$ \final{98.3\%}, $\approx$ \final{745 min (12 h) a month}. (b) Each stateless service becomes $1 - 0.005^3 \approx 99.99999\%$: \final{99.79\%} ($\approx$ 92 min). (c) Multi-AZ SQL: \final{99.84\%} ($\approx$ 70 min); second provider ($1 - 0.001^2$): \final{99.94\%} ($\approx$ 26 min). The unreplicated links dominate, in order.
\careful{replicas whose failures are correlated (the same bad deploy, the same zone): then $1 - (1-a)^n$ is fiction.}
\ans{5} (a) $100 \times 0.2 =$ \final{20}; $100 \times 1 =$ \final{100}; $100 \times 4 =$ \final{400 $>$ 200}: \texttt{order} runs out of threads, the gateway waiting on \texttt{order} runs out too, and every route behind them fails (\final{cascade}). (b) \final{At most 100}. (c) $50/20 =$ \final{2.5$\times$} (latency up to $\approx$ 500 ms); at 4 s only 50 threads wait and \final{150 keep serving}. (d) $2 \times 2 + 0.3 \approx$ \final{4.3 s $\le$ 5 s $\le$ 6 s: it fits}.
\careful{leaving the client's default timeout (often minutes) or setting a timeout longer than the caller's own deadline.}
\ans{6} (a) $3 \times 3 \times 3 =$ \final{27 calls}: retry in \final{one layer only}, 1--2 retries. (b) \final{200, 400 and 800 ms}; with full jitter the third wait is \final{random between 0 and 800 ms}, so clients do not come back together. (c) Connect timeout: \final{retry}. Read timeout without a key: \final{never} (the charge may have happened; with an idempotency key, retry). 503 with Retry-After: \final{retry after the indicated wait}. 409 and card declined: \final{never} (business answers). (d) \final{A fast, honest error} (``payment temporarily unavailable, your seats are held for 10 minutes''), not a 4 s wait; after 30 s three probes test the provider.
\careful{counting business errors (card declined) as failures, which opens the breaker for everybody.}
\ans{7} (a) \final{No}: the standby replicates the \texttt{DELETE} within milliseconds. \final{A backup with point-in-time recovery} (restore to 21:59:30) saves you; replicas are not backups. (b) $2 \times 10^{12}/2 \times 10^8 = 10\,000$ s $\approx$ \final{2.8 h}, far above 15 min: a \final{warm standby} (or a delayed replica) in another region. (c) $2 + 5 + 1 + 5 + 1 =$ \final{14 min}; with a 1-hour TTL, \final{part of the traffic keeps going to the dead region for up to an hour}. (d) Two regions behind DNS with health checks, region B with 2 replicas ($\approx$ 10\%) ready to scale. Experiment: \final{hypothesis} ``with the provider at 3 s the breaker opens within 30 s and catalog p99 stays $<$ 300 ms''; \final{fault} 3 s $\pm$ 0.5 s of latency (Toxiproxy) on one canary; \final{abort} if catalog p99 $>$ 1 s or burn $>$ 10$\times$.
\careful{never testing a restore: an untested backup is not a backup.}
